%---------------------------------------------------------------------%
\addChapter{Lampiran 1 - Kode Program}
\chapter*{Lampiran 1 - Kode Program}
%---------------------------------------------------------------------%

\section*{Lampiran 1.1 Contoh Penulisan Kode Program (\texttt{contoh.py})}
\addcontentsline{toc}{section}{Lampiran 1.1 Contoh Penulisan Kode Program (contoh.py)}

Lampiran digunakan untuk memuat berkas pendukung seperti kode program, data
mentah, atau dokumen tambahan. Potongan kode ditulis menggunakan lingkungan
\texttt{lstlisting} dengan gaya \texttt{mystyle} yang telah didefinisikan pada
\texttt{thesis.tex}.

\begin{lstlisting}[
    style=mystyle,
    language=Python,
]
def main():
    """Ganti isi lampiran ini dengan kode program yang sebenarnya."""
    data = [1, 2, 3, 4, 5]
    total = sum(data)
    print("Jumlah data:", len(data))
    print("Total nilai:", total)


if __name__ == "__main__":
    main()
\end{lstlisting}

%---------------------------------------------------------------------%
\section*{Lampiran 1.2 Contoh Lampiran Tabel}
\addcontentsline{toc}{section}{Lampiran 1.2 Contoh Lampiran Tabel}

Tabel pendukung yang terlalu panjang untuk dimuat pada batang tubuh laporan
dapat diletakkan pada lampiran.

\begin{table}[htbp]
    \centering
    \addtocounter{table}{-1}
    \caption{Contoh tabel pada lampiran}
    \label{tab:contoh_lampiran}
    \small
    \setlength{\tabcolsep}{4pt}
    \renewcommand{\arraystretch}{1.15}
    \begin{tabularx}{\textwidth}{
        >{\centering\arraybackslash}p{2cm}
        >{\raggedright\arraybackslash}p{3cm}
        >{\raggedright\arraybackslash}X
    }
        \toprule
        \textbf{No.} & \textbf{Parameter} & \textbf{Keterangan} \\
        \midrule
        1 & Parameter A & Keterangan parameter pertama. \\
        2 & Parameter B & Keterangan parameter kedua. \\
        3 & Parameter C & Keterangan parameter ketiga. \\
        \bottomrule
    \end{tabularx}
\end{table}
